\documentclass[addpoints]{exam}
\usepackage{amsmath,amssymb,graphicx}
\usepackage[colorlinks]{hyperref}
\usepackage{pgfplots}
\usepgfplotslibrary{units}
\pgfplotsset{compat=1.9}
\runningheadrule
\runningheader{}{CS 378H Spring 2023 Linear Algebra Worksheet}{}
\runningfooter{}{Page \thepage}{}
\newcommand{\bv}{\mathbf{v}}
\newcommand{\bu}{\mathbf{u}}
\newcommand{\bw}{\mathbf{w}}
\newcommand{\bp}{\mathbf{p}}
\newcommand{\bc}{\mathbf{c}}
\begin{document}
\title{CS 378H Spring 2024 Linear Algebra Worksheet}
\author{Due Tuesday, January 23, 2024 at 1:59 PM}
\date{\parbox{\textwidth}{\small ``There is hardly any theory which is more
elementary [than linear algebra], in spite of the fact that generations of
professors and textbook writers have obscured its simplicity by preposterous
calculations with matrices.'' --J.~Dieudonne}}
\maketitle
\vspace{0.1in}
\noindent
\makebox[\textwidth]{Name: Arham Lodha}
\vspace{0.2in}
\noindent
\makebox[\textwidth]{UT EID: arl3759}
\vspace{0.2in}
\begin{center}
\fbox{\fbox{\parbox{5.5in}{
Questions begin on the next page. Please answer each question in the spaces
provided; if you need extra room, you can attach extra blank pages at the end of
the worksheet.
\vspace{0.1in}
The worksheet follows the standard class collaboration policy. You may ask other
students, the instructor, or the TA for help, but \textbf{all work written on this
worksheet must be entirely your own.} Do not simply copy somebody else's answers.
\vspace{0.1in}
You may look at linear algebra textbooks, online tutorials, math.stackexchange, or
any other resources to help you complete the worksheet. Do not worry if you cannot
do the problems from memory alone! You are allowed, and expected, to look up any
formulas or definitions as needed.
\vspace{0.1in}
The problems are not intended to be difficult or time-consuming to solve, but
rather, to refresh your memory on basic linear algebra concepts, and to give you a
self-check on the prerequisites you will need to succeed in this class. If you get
stuck on any problem, the TA or instructor would be happy to help you during office
hours.
}}}
\end{center}
\vspace{0.5in}
\section*{Useful Formulas}
\begin{itemize}
\item Dot product: $(u_x,u_y,u_z,u_w) \cdot (v_x,v_y,v_z,v_w) = u_xv_x + u_yv_y +
u_zv_z + u_wv_w$.
\item Cross product: $(u_x,u_y,u_z) \times (v_x,v_y,v_z) = (u_y v_z - u_z v_y, u_z
v_x - u_x v_z, u_x v_y - u_yv_x)$.
\item Norm: $\|\bu\| = \sqrt{\bu \cdot \bu} = \sqrt{\bu^T \bu}$.
\item Unit vector: $\hat{\bu} = \frac{\bu}{\|\bu\|}.$
\item Angles between vectors: $\bu \cdot \bv = \|\bu\|\|\bv\|\cos \theta, \quad \|\
vec \bu \times \vec \bv\| = \|\bu\|\|\bv\||\sin \theta|.$
\end{itemize}
\newpage
\begin{questions}
    \question
    Recall that a function $f(\bv)$ of a vector $\bv$ is \emph{linear} if $f(\alpha \
    bv) = \alpha f(\bv)$ and $f(\bv+\bw) = f(\bv) + f(\bw).$ For each of the following
    functions, (a) state whether or not the function is linear, and (b) if the function
    is not linear, give a counterexample where the function violates one of the above
    properties. You do \emph{not} need to provide a proof if the function is linear;
    siimply stating it's linear is enough.
    \begin{parts}
    \part[1] Adding a non-zero constant vector $\bc$ to the input vector, $f(\bv) = \
    bv+\bc$.

    a) The function is not linear

    b) $\forall v \in V$, $f(0v) = f(0) = b(0) + c = c \neq 0f(v) = 0(bv + c) = 0$  


    \part[1] The constant function $f(\bv) = 42.$
    
    a) the function is not linear

    b) $\forall v \in V$, $f(0) = 42 \neq 0f(v) = 0$  

    \part[1] Swapping the coordinates of a 2D vector: $f(\bv) = (\bv_y, \bv_x)$
    
    a) The function is linear

    \part[1] Vector norm, $f(\bv) = \|\bv\|$.
    
    a) The function is not linear.

    b) Let $v = \begin{bmatrix}
        1 \\ 0
    \end{bmatrix}$ and let $w = \begin{bmatrix}
        0 \\ 1
    \end{bmatrix}$. $f(v) = 1 = f(w)$. $f(v + w) = \sqrt{2} \neq f(v) + f(w) = 2$  

    \part[2] The function returning the determinant of a $2\times 2$ matrix: $f(M) = \
    det M$.

    a) Not linear

    b) $f(2I) = \det(2I) = 4 \neq 2f(I) = 2$  

    \part[2] Cross product with a fixed vector $\bc$: $f(\bv) = \bv \times \bc$.

    a) Function is linear

    \part[2] Angle $\theta(\bv) = \angle(\bv,\bc)$ between $\bv$ and a fixed vector $\bc$.

    a) Function is not linear

    b) $v = \begin{bmatrix}
        1 \\ 0
    \end{bmatrix}$, $c = w = \begin{bmatrix}
        0 \\ 1
    \end{bmatrix}$, $\theta(v) = \frac{\pi}{2}$, $\theta(w) = 0$. $\theta(v + w) = \frac{\pi}{4} \neq \theta(v) + \theta(w) = \frac{\pi}{2}$     

    \end{parts}

    \newpage
\question
Although in practice you will not need to do a lot of linear algebra by hand
(that's what computers are for!), you should be able to perform basic computations,
and know how to use computational resources like Matlab, Mathematica, Wolfram Alpha
to help you debug your code.
\begin{parts}
\part[3] A rotation matrix $R$ is a square matrix satisfying the equations $R^TR =
I$ and $\det R = 1$, where $I$ is the identity matrix. Is the following matrix
(approximately) a rotation marix? How do you know? \textbf{You do not need to do
the calculation by hand or show work.}
$$\left[\begin{array}{cc} 0.743 & -0.669 \\ 0.669 & 0.743 \end{array}\right]$$

Yes, it is approximately a rotation matrix because $$\begin{bmatrix}
    0.743 & -0.669 \\ 0.669 & 0.743
\end{bmatrix}^T \begin{bmatrix}
    0.743 & -0.669 \\ 0.669 & 0.743
\end{bmatrix} = \begin{bmatrix}
    99961/100000 & 0\\
           0 & 99961/100000\\
\end{bmatrix} \approx I$$. 


\part[3] Does the matrix
$$M = \left[\begin{array}{ccc}1 & 2 & 3\\4 & 5 & 6\\7 & 8 & 9\end{array}\right]$$
have an inverse? How do you know for sure? If an inverse exists, compute it. \textbf{You do not need to do the calculation by hand or show work.}

No, M is not invertible because $\det(M) = 0$. 

\part[4] Use your favorite computer linear algebra tool to solve
$$\left[\begin{array}{cccc}2 & 1 & 7 & 8\\-1 & 1 & 2 & -3\\1 & 3 & 5 & 7\\ 0 & 0 &
0 & 1\end{array}\right]\bv = \left[\begin{array}{c}1\\2\\3\\4\end{array}\right]$$
for $\bv$. Round each coordinate of your answer to the nearest 0.01. \textbf{You do
not need to do the calculation by hand or show work.}

$v = \begin{bmatrix}
    -16.52 \\ -4.35 \\ 0.91 \\ 4
\end{bmatrix}$ 

\end{parts}
\newpage


\question
Let $\bc\in \mathbb{R}^2$ be a fixed vector. Let
$$f(\bv) = \left(I - 2\frac{\bc\bc^T}{\bc^T\bc}\right)\bv,$$
for vectors $\bv\in\mathbb{R}^2$ in the plane, where $I = \left[\begin{array}{cc}1
& 0\\0 & 1\end{array}\right]$ is the identity matrix.
\begin{parts}
\part[1] Compute $\bc\bc^T$ and $\bc^T \bc$ for $\bc = \begin{bmatrix}1 \\ -2\end{bmatrix}$. Are these similar-looking expressions actually the same, or
different?

$\bc \bc^T = \begin{bmatrix}
    1 \\ -2
\end{bmatrix} \begin{bmatrix}
    1 & -2
\end{bmatrix} = \begin{bmatrix}
    1 & -2 \\ -2 & 4
\end{bmatrix}$

$\bc^T \bc = \begin{bmatrix}
    1 & -2
\end{bmatrix} \begin{bmatrix}
    1 \\ -2
\end{bmatrix} = 1 + 4 = 5$

No, the similar looking expressions are not the same, for $c = \begin{bmatrix}
    1 & -2
\end{bmatrix}$ the first expression is a $2 \times 2$ matrix while the second is a scalar.     

\part[1] Is the function $f(\bv)$ linear? If not, provide a counterexample.

Yes it is linear. $\forall v, w \in \mathbb{R}^2$ and $\forall a, b \in \mathbb{R}$, $f(av + bw) = \left(I - 2\frac{\bc\bc^T}{\bc^T\bc}\right)(av + bw) = a\left(I - 2\frac{\bc\bc^T}{\bc^T\bc}\right)v + b\left(I - 2\frac{\bc\bc^T}{\bc^T\bc}\right)w = af(v) + bf(w)$. 

\part[2] Compute $f(\bc)$.

$f(\bc) = \left(I - 2\frac{\bc\bc^T}{\bc^T\bc}\right)\bc = c - 2\frac{\bc\bc^T\bc}{\bc^T\bc} = -c$ 

\part[2] Compute $f(\bv)$, where $\bv$ is any vector perpendicular to $\bc$ (i.e.
$\bv \cdot \bc = 0$).

$f(v) = \left(I - 2\frac{\bc\bc^T}{\bc^T\bc}\right)v = v - 2\frac{\bc\bc^Tv}{\bc^T\bc} = v - 2\frac{\bc (c \cdot v)}{\bc^T\bc} = v - 2(0) = v$ 

\part[4] Describe using words, and no equations, what the function $f$ does to
vectors. You may find it helpful to probe the behavior of $f(\bv)$ for concrete
fixed vectors $\bc$ and draw some pictures. Your answers to parts (c) and (d)
should also prove useful.
\end{parts}
From the input vector $v \in \mathbb{R}^2$, f subtracts twice the portion of v which is in the direction of $c$, in other words it inverts the portion of v which is in the direction of the fixed vector $c$. More simply, it reflects the input vector v with respect to the plane normal to $c$.       

\question
A key skill in linear algebra is turning systems of linear equations into matrix
equations of the form $A\mathbf{x} = \mathbf{b}$, where $A$ and $\mathbf{b}$ are
known, and $\mathbf{x}$ is unknown.
\begin{parts}
\part[5] Let $\bv_1,\bv_2,\bv_3$ and $\bw$ be four different vectors in $\mathbb{R}^3$. You are told that
$$\bw = \alpha \bv_1 + \beta \bv_2 + \gamma \bv_3$$
for some unknown real numbers $\alpha, \beta$, and $\gamma$. Write a matrix
equation (of the form $A\mathbf{x} = \mathbf{b}$) that would allow you to solve for
these unknown scalars. (\textbf{You do not need to actually solve the equation.})

$A = \begin{bmatrix}
    \vline & \vline & \vline \\
    \bv_1 & \bv_2 & \bv_3 \\ 
    \vline & \vline & \vline \\
\end{bmatrix} = \begin{bmatrix}
    v_{1_x} & v_{2_x} & v_{3_x} \\
    v_{1_y} & v_{2_y} & v_{3_y}\\
    v_{1_z} & v_{2_z} & v_{3_z}\\
\end{bmatrix}$ 

$x = \begin{bmatrix}
    \alpha \\ \beta \\ \gamma
\end{bmatrix}$ 

$Ax = \begin{bmatrix}
    \vline & \vline & \vline \\
    \bv_1 & \bv_2 & \bv_3 \\ 
    \vline & \vline & \vline \\
\end{bmatrix} \begin{bmatrix}
    \alpha \\ \beta \\ \gamma
\end{bmatrix} = \begin{bmatrix}
    v_{1_x} & v_{2_x} & v_{3_x} \\
    v_{1_y} & v_{2_y} & v_{3_y}\\
    v_{1_z} & v_{2_z} & v_{3_z}\\
\end{bmatrix}\begin{bmatrix}
    \alpha \\ \beta \\ \gamma
\end{bmatrix} = w$

\part[5] Let $\bv_1,\bv_2,\bv_3$ be three different vectors in $\mathbb{R}^3$, and
$k_1, k_2, k_3$ three real numbers. You are told that
\begin{align*}
k_1 &= \bw \cdot \bv_1\\
k_2 &= \bw \cdot \bv_2\\
k_3 &= \bw \cdot \bv_3
\end{align*}
for some unknown vector $\bw\in\mathbb{R}^3$. Write down a matrix equation that
would allow you to solve for this vector $\bw$.

$A = \begin{bmatrix}
    v_{1_x} & v_{1_y} & v_{1_z} \\
    v_{2_x} & v_{2_y} & v_{2_z}\\
    v_{3_x} & v_{3_y} & v_{3_z}\\
\end{bmatrix}$ 

$x = w$

$b = \begin{bmatrix}
    k_1 \\ k_2 \\ k_3
\end{bmatrix}$ 

$Aw = \begin{bmatrix}
    v_{1_x} & v_{1_y} & v_{1_z} \\
    v_{2_x} & v_{2_y} & v_{2_z}\\
    v_{3_x} & v_{3_y} & v_{3_z}\\
\end{bmatrix}w = b$ 



\end{parts}

\newpage
\question
Suppose you wrote a computer program that drew an animation by plotting the motion
of a point with coordinates $\bp(t)$ in the plane, as a function of time $t\in [0,\infty).$ Describe what you would see on the screen, as time evolves. (The intent is
not for you to actually write such a program, but doing so is not forbidden if
you're stuck.) If you sketch pictures, clearly indicate how the figure is changing
over time.
\begin{parts}
\part[5] $$\bp(t) = (1-t)\left[\begin{array}{c}-3\\4\end{array}\right] + t\left[\begin{array}{c}-3\\5\end{array}\right].$$

A point starting at $\begin{bmatrix}
    -3 & 4
\end{bmatrix}$ and going in the $\begin{bmatrix}
    0 \\ 1
\end{bmatrix}$ direction in a straight line.   

\part[5] $$\bp(t) = \left[\begin{array}{cc} \cos(t^2) & -\sin(t^2)\\ \sin(t^2) & \cos(t^2)\end{array}\right]\left[\begin{array}{c}e^{t}\\0\end{array}\right].$$

A point starting at $\begin{bmatrix}
    1 \\ 0
\end{bmatrix}$ and going around a spiral whose points are getting exponentially further from the origin with respect to time time .

\end{parts}

\end{questions}
\end{document}